% !TeX root = ../../../金融经济学_王江_讲义及习题.tex
\begin{tikzpicture}[>=Stealth,font=\footnotesize,x=1.16cm,y=1.23cm,
    curve/.style={draw=blue!70!black,very thick},
    chord/.style={draw=red!65!black,thick,densely dashed}]
 \foreach \panelshift/\curvesign/\strict/\curvelevel/\panelname in {
   0/1/0/.88/{(a) 凸（非严格）},
   3.55/1/1/1.03/{(b) 严格凸},
   7.10/-1/0/1.13/{(c) 凹（非严格）},
   10.65/-1/1/.98/{(d) 严格凹}
 } {
   \begin{scope}[xshift=\panelshift cm]
     \draw[->] (0,0)--(2.60,0) node[right,inner sep=1pt] {$x$};
     \draw[->] (0,0)--(0,1.95) node[above] {$f(x)$};
     \ifnum\strict=1
       \draw[curve,domain=.20:2.26,samples=45]
         plot (\x,{\curvelevel+\curvesign*.76*(\x-1.23)^2});
       \pgfmathsetmacro{\chordheight}{\curvelevel+\curvesign*.76*.65^2}
     \else
       \draw[curve,domain=.20:2.26,samples=55]
         plot (\x,{\curvelevel+\curvesign*13*max(abs(\x-1.23)-.69,0)^3});
       \let\chordheight\curvelevel
     \fi
     \draw[chord] (.58,\chordheight)--(1.88,\chordheight);
     \fill[fill=red!65!black] (.58,\chordheight) circle(1.35pt)
       (1.88,\chordheight) circle(1.35pt);
     \node at (1.23,-.32) {\panelname};
   \end{scope}
 }
\end{tikzpicture}
